% CHAPTER 2 - LITERATURE REVIEW
\chapter{Literature Review}

\section{Introduction}
This chapter reviews the research that the Context Memory Framework builds on. The project
sits at the intersection of four lines of work. The first is \textit{Natural Language to
SQL}, which has moved from task-specific neural parsers to prompting general-purpose Large
Language Models. The second is \textit{retrieval}, in particular retrieval-augmented
generation, dense and sparse retrieval, rank fusion and reranking. The third is
\textit{memory for LLM-based agents}, which organizes what a system remembers across
interactions. The fourth is the \textit{security} of LLM-integrated applications, where
prompt injection and unsafe tool use are now well documented.

Each of these areas solves part of the problem described in Chapter~1. NL2SQL research
shows how to produce SQL from a question and a schema; retrieval research shows how to
select relevant text from a large store; memory research shows how to keep and organize
knowledge over time; and security research shows what can go wrong when model output is
executed. The final section identifies what is missing when these pieces are used in
isolation, which defines the contribution of this project.

\section{Evolution of Natural Language to SQL}

\subsection{Early Neural Approaches and Benchmarks}
Translating questions into database queries is a long-standing problem in natural
language interfaces to databases. Neural sequence-to-sequence models made the task
data-driven. Seq2SQL \cite{zhong2017seq2sql} introduced the WikiSQL dataset and generated
single-table queries, using reinforcement learning with query execution as a reward signal
for the unordered parts of the query.

Spider \cite{yu2018spider} raised the difficulty by introducing a large, cross-domain
benchmark of 10,181 questions and 5,693 unique complex SQL queries over 200 multi-table
databases covering 138 domains, with databases in the test set unseen during training. This
forced models to generalize to new schemas, joins, nesting and grouping rather than
memorize a fixed database. BIRD \cite{li2023bird} moved closer to real enterprise use with
12,751 question--SQL pairs over 95 large databases totalling 33.4\,GB across 37
professional domains. BIRD explicitly provides \textit{external knowledge evidence} for
many questions, recognizing that the meaning of a question frequently depends on domain
knowledge that cannot be read from the schema.

\subsection{Schema Linking}
A central sub-problem is \textit{schema linking}: deciding which tables and columns a
question refers to. RESDSQL \cite{li2023resdsql} decouples schema linking from skeleton
parsing and uses a ranking-enhanced encoder that keeps only the most relevant schema items
in the model input. This idea, selecting a relevant subset of the schema before
generation, is important for large schemas and is used in this project through vector
search over table descriptions combined with foreign-key expansion.

\section{Large Language Models for Text-to-SQL}
With the arrival of large code-capable language models, NL2SQL moved from training
specialized parsers to prompting. Rajkumar et al.\ \cite{rajkumar2022evaluating} evaluated
Codex on Spider and found that a prompt containing the database's \texttt{CREATE TABLE}
statements and a few sample rows gave competitive results without fine-tuning. DIN-SQL
\cite{pourreza2023dinsql} decomposed the task into schema linking, query classification and
decomposition, generation, and self-correction, showing that structuring the prompt
pipeline improves accuracy. DAIL-SQL \cite{gao2023dailsql} systematically compared question
representations, example selection and example organization for in-context learning, and
showed that retrieving \textit{similar past question--SQL pairs} as examples is an effective
form of context.

These results share one message: with capable models, accuracy depends strongly on
\textit{what} context is placed in the prompt and \textit{how} it is organized. They also
share an assumption that the needed knowledge is contained in the schema, the question, the
examples, or (in BIRD) an evidence string supplied with the benchmark. In an organization
nobody supplies that evidence per question; it has to be remembered and retrieved.

\begin{table}[H]
    \centering
    \caption{Representative NL2SQL Work and the Context Each Approach Relies On}
    \begin{tabular}{p{3.2cm} p{4.6cm} p{5.2cm}}
        \toprule
        \textbf{Work} & \textbf{Main idea} & \textbf{Context used for generation} \\
        \midrule
        Seq2SQL \cite{zhong2017seq2sql} & Seq2seq with RL on execution & Question + single table \\[4pt]
        Spider \cite{yu2018spider}      & Cross-domain benchmark & Question + full schema \\[4pt]
        RESDSQL \cite{li2023resdsql}    & Ranked schema filtering & Question + top schema items \\[4pt]
        BIRD \cite{li2023bird}          & Large realistic databases & Question + schema + given evidence \\[4pt]
        DIN-SQL \cite{pourreza2023dinsql} & Decomposed prompting & Question + schema links + sub-steps \\[4pt]
        DAIL-SQL \cite{gao2023dailsql}  & Example selection study & Question + schema + similar examples \\[4pt]
        \textbf{CMF (this work)}        & Governed memory layer & Question + selected schema + approved memories + examples + history \\
        \bottomrule
    \end{tabular}
    \label{tab:nl2sql_survey}
\end{table}

\section{Retrieval-Augmented Generation and Hybrid Search}

\subsection{Retrieval-Augmented Generation}
Lewis et al.\ \cite{lewis2020rag} introduced Retrieval-Augmented Generation (RAG), which
combines a parametric generator with a non-parametric memory: a retriever selects passages
from an external index and the generator conditions on them. RAG makes knowledge updatable
without retraining and lets the system show which sources were used. CMF applies the same
pattern, but the ``documents'' are short business memories, table descriptions and past
question--SQL pairs rather than long passages.

\subsection{Sparse and Dense Retrieval}
Classical retrieval scores documents by term statistics. The BM25 function
\cite{robertson2009bm25} weights each matching term by its inverse document frequency and a
saturating function of its frequency in the document, so that repeated terms contribute
with diminishing returns. Sparse methods are precise for exact identifiers such as column
names (\texttt{invoice\_date}), which matter greatly in NL2SQL.

Dense retrieval encodes queries and documents as vectors in a shared space. Dense Passage
Retrieval \cite{karpukhin2020dpr} showed that a learned dual encoder can outperform BM25 for
open-domain question answering. Sentence-BERT \cite{reimers2019sbert} produced sentence
embeddings suitable for cosine-similarity search, and MiniLM \cite{wang2020minilm}
distillation produced compact models; the \texttt{all-MiniLM-L6-v2} model used in this
project produces 384-dimensional embeddings. Dense vectors capture paraphrases such as
``uncollectible receivables'' versus ``write-offs'', but can miss exact tokens. Approximate
nearest-neighbour indexes such as HNSW \cite{malkov2018hnsw} make dense search fast at scale
and are used by vector databases such as Qdrant \cite{qdrantdocs}.

\subsection{Rank Fusion and Reranking}
Because sparse and dense retrieval fail in different ways, they are often combined.
Reciprocal Rank Fusion (RRF) \cite{cormack2009rrf} merges ranked lists without needing
comparable scores:
\begin{equation}
    \text{RRF}(d) = \sum_{r \in \mathcal{R}} \frac{1}{k + \text{rank}_r(d)}
\end{equation}
where $\mathcal{R}$ is the set of rankings and $k$ is a constant (60 in the original
paper). Rank fusion is simple and robust, which is why it is offered natively by modern
vector databases.

A second stage then reorders a small candidate set with a more expensive model. Nogueira
and Cho \cite{nogueira2019passage} showed that a BERT cross-encoder, which reads the query
and the passage together, substantially improves passage ranking on MS MARCO
\cite{bajaj2016msmarco}. The \texttt{cross-encoder/ms-marco-MiniLM-L-6-v2} reranker used in
this project follows this design.

\section{Memory Architectures for LLM Agents}

\subsection{Cognitive Memory Categories}
Cognitive psychology distinguishes \textit{semantic} memory (general facts and meanings)
from \textit{episodic} memory (specific experienced events) \cite{tulving1972episodic}, and
both from \textit{procedural} memory (how to do things). The CoALA framework
\cite{sumers2023coala} applies this taxonomy to language agents, describing a short-term
\textit{working memory} and long-term episodic, semantic and procedural memories, together
with actions that retrieve from and write to them. CMF adopts this vocabulary directly:
business definitions are semantic, formatting and column-usage preferences are procedural,
notable analytical events are episodic, and the current conversation plus the last
structured query form working memory.

\subsection{Retrieval from Long-Term Memory}
Generative Agents \cite{park2023generative} store an agent's experiences in a memory stream
and retrieve them by a combination of \textit{recency} (exponential decay), \textit{importance}
(a score assigned when the memory is written) and \textit{relevance} (embedding similarity).
MemGPT \cite{packer2023memgpt} treats the limited context window like main memory in an
operating system and pages information in and out of external storage. Both show that what
reaches the prompt should be selected by an explicit policy rather than by appending
history.

CMF's first-stage score follows the same principle, extending it with keyword overlap and
a memory-type match, and adding a cross-encoder reranker. Unlike these agent systems, CMF
operates in a multi-user organization, so memory retrieval must also respect ownership,
sharing and database scope, and memories written by the model must be reviewed by people.

\section{Context Engineering and Long Prompts}
Longer context windows do not guarantee that information is used. Liu et al.\
\cite{liu2023lost} showed that models use information at the beginning and end of a long
context more reliably than information in the middle (``lost in the middle''). This
supports selecting a small, relevant context and placing the most important material in
deliberate positions. In CMF, the system rules and the live schema always come first, the
current question and the generation requirements always come last, and optional sections
are admitted in priority order under a token budget.

A second concern is the \textit{authority} of each part of the prompt. When retrieved text
is mixed with instructions, the model cannot reliably tell them apart. Making the priority
order explicit, and labelling retrieved memories as untrusted data, is a lightweight form of
the separation that security research recommends.

\section{Security of LLM-Generated SQL}
Greshake et al.\ \cite{greshake2023indirect} demonstrated \textit{indirect prompt
injection}: instructions planted in data that an LLM-integrated application later
retrieves can take control of the model's behaviour. Pedro et al.\ \cite{pedro2023prompt}
studied LLM-integrated web applications built with LangChain and showed
\textit{prompt-to-SQL} injection attacks in which a user's natural-language input leads the
model to generate destructive or data-leaking SQL. The OWASP Top~10 for LLM Applications
\cite{owasp2023llm} lists prompt injection and excessive agency among the main risks of
such systems.

For NL2SQL, these findings mean that (i) stored memories are a persistent injection
surface, (ii) instructions in the system prompt are not a sufficient defence, and (iii) the
application should limit what generated statements can do independently of the model.
Parsing SQL into an abstract syntax tree, as the \texttt{sqlglot} library does
\cite{sqlglotdocs}, allows a deterministic classification of each statement before
execution. Human approval for any change that modifies data limits the agency of the
model to proposing, not executing.

\section{Literature Gap}
The reviewed work provides strong individual components but leaves a gap for NL2SQL
inside organizations:
\begin{itemize}[leftmargin=*, topsep=4pt, itemsep=4pt]
    \item NL2SQL research assumes that external knowledge is given with each question
          \cite{li2023bird}; it does not address how such knowledge is captured from
          conversations, reviewed, versioned and kept current.
    \item Agent memory systems \cite{park2023generative, packer2023memgpt} rank memories by
          relevance, recency and importance but do not combine ranking with multi-tenant
          authorization, human approval and explicit superseding of outdated rules.
    \item Hybrid retrieval and reranking \cite{cormack2009rrf, nogueira2019passage} are
          rarely applied jointly to business memories, schema tables and past question
          examples in one pipeline whose decisions are visible to the user.
    \item Security work describes prompt and SQL injection risks
          \cite{greshake2023indirect, pedro2023prompt} but few NL2SQL prototypes integrate
          a deterministic statement planner, read-only execution, masking and approved,
          re-validated data changes.
\end{itemize}
The Context Memory Framework addresses this gap with a single, inspectable pipeline that
combines governed long-term memory, authorization-first hybrid retrieval, prioritized
context assembly and deterministic SQL guardrails.
